%% Motivic Shadow Tower on the Class-M Depth-1 Lane -- Platonic Chapter.
%%
%% Companion chapter to chapters/theory/motivic_shadow_tower.tex, which
%% established motivic rationality of the Virasoro shadow tower
%% (thm:virasoro-motivic-rationality-all-r). This file records the
%% conditional class-M extension on the depth-1 residue lane:
%%
%% (1) W_N principal on the conditional depth-1 generator-line lane;
%% (2) Bershadsky-Polyakov T-line in Arakawa convention;
%% (3) W_infinity at generic Psi on the conditional T-line lane;
%% (4) Tempered Virasoro (remark: not independently realized);
%%
%% and isolates the structural reason on this lane: shadow-tower
%% residues on Conf_n(D) extract depth-1 Arnold data, hence rational
%% outputs. The E_2-topological Drinfeld associator, by contrast,
%% carries depth-2 MZV via iterated Arnold integrals; this is the sharp
%% E_1-chiral vs E_2-topological distinction after the coefficient's
%% depth has been fixed.
%%
%% ARCHITECTURE. Six self-contained inscriptions:
%%
%% Section 1. Structural theorem: shadow residues are depth-1 data.
%% (thm:shadow-tower-depth-1-rationality,
%% thm:e1-vs-e2-mzv-depth-distinction)
%%
%% Section 2. Conditional W_N principal rationality.
%% (thm:w-n-motivic-rationality-all-r,
%% prop:w3-w-line-motivic-rationality)
%%
%% Section 3. Bershadsky-Polyakov T-line rationality.
%% (thm:bp-motivic-rationality-arakawa,
%% prop:bp-fl-convention-caveat)
%%
%% Section 4. Conditional W_infinity rationality at generic Psi.
%% (thm:w-infty-motivic-rationality-all-r)
%%
%% Section 5. Conditional class-M depth-1 theorem.
%% (thm:class-m-motivic-rationality-full)
%%
%% Section 6. Tempered Virasoro remark and frontier.
%% (rem:tempered-virasoro-not-realized)

\providecommand{\PhiFA}{\Phi^{\mathrm{FA}}}
\providecommand{\SpCh}{\mathrm{Sp}^{\mathrm{ch}}}

\chapter{Motivic rationality on the class-M depth-1 lane}
\label{ch:motivic-class-m-full}

\begin{overture}
The Virasoro shadow tower is motivic-rational
(Theorem~\ref{thm:virasoro-motivic-rationality-all-r}); no
$\zet^{\mot}$ enters at any weight. This does not imply that every
class-M shadow coefficient is rational. The obstruction is local:
when the coefficient is represented by a finite OPE-residue polynomial
on the $E_1$-chiral bar side, the computation sees depth-$1$ Arnold
data. When the coefficient uses an $E_2$ comparison or an iterated
configuration-space integral, the Drinfeld associator may enter. The
chapter records the proved Virasoro and BP T-line cases and the
conditional principal $\cW_N$, $\cW_{\infty}$, and class-M extensions
on this depth-$1$ lane.
\end{overture}

%%%%%%%%%%%%%%%%%%%%%%%%%%%%%%%%%%%%%%%%%%%%%%%%%%%%%%%%%%%%%%%%%%%%%%%%%%%%%
\section{Depth-1 rationality of the shadow-tower residue}
\label{sec:depth-1-rationality}
%%%%%%%%%%%%%%%%%%%%%%%%%%%%%%%%%%%%%%%%%%%%%%%%%%%%%%%%%%%%%%%%%%%%%%%%%%%%%

The shadow coefficient $S_r(\cA)$ is extracted from the universal
Maurer--Cartan element $\Theta_{\cA}$ on the bar complex
$B^{\ord}(\cA) = T^c(s^{-1} \bar{\cA})$ on $\mathrm{Conf}_n(D)$.
When this extraction is represented by a finite residue polynomial in
OPE coefficients, it sees depth-$1$ Arnold data and not iterated
Arnold periods. Without this hypothesis, higher configuration-space
periods can enter.

\begin{theorem}[\label{thm:shadow-tower-depth-1-rationality}Shadow
residue rationality, depth-$1$ lane]
\ClaimStatusConditional
Let $\cA$ be a chirally Koszul algebra on the formal disk $D$ with
$\mathbb{Q}$-rational OPE structure constants. Fix $r \geq 2$ and
assume that $S_r(\cA)$ is computed by a finite residue expression in
the OPE coefficients, with no iterated Arnold integral and no
associator-dependent comparison term. Then $S_r(\cA)$ is a
$\mathbb{Q}$-rational function of the family parameter and its motivic
lift has weight zero:
\[
 S_r^{\mot}(\cA) \in
 \MZV^{\mot}_0 \otimes \mathbb{Q}(\mathrm{param})
 =
 \mathbb{Q}(\mathrm{param}),
\]
so no positive-weight $\zet^{\mot}(n)$ or $(2\pi i)^{\mot,2n}$
appears on this lane.
\end{theorem}

\begin{proof}
Let $n \geq 2$. The $n$-th bar term
$B_n^{\ord}(\cA) = (s^{-1} \bar{\cA})^{\otimes n}$ carries the Arnold
propagator $\omega_{ij} = d\log(z_i - z_j)$ on $\mathrm{Conf}_n(D)$,
a $1$-form of motivic weight $1$ (Brown 2012 Thm.\ 1.1: logarithmic
$1$-forms generate $\MZV^{\mot}_1$, and $\mathbb{Q}$-linear
combinations of residues of logarithmic forms lie in $\MZV^{\mot}_0 =
\mathbb{Q}$). The shadow coefficient extraction is
\[
 S_r(\cA) \;=\; \operatorname{Res}_{z_1, \ldots, z_n}\!
 \Bigl[\,
 \Theta_{\cA}(z_1, \ldots, z_n) \cdot \prod_{i < j} \omega_{ij}
 \,\Bigr],
\]
a sum of \emph{simple residues} of products of Arnold forms against
$\Theta_{\cA}$. Each simple residue is a $\mathbb{Q}$-linear functional
on the OPE residues of $\cA$; under the assumption that the OPE
residues lie in $\mathbb{Q}(\mathrm{param})$, the sum lies in
$\mathbb{Q}(\mathrm{param})$.

\emph{Depth-$1$ restriction.} An iterated Arnold integral
$\int\!\int \omega_{ij} \wedge \omega_{kl}$ of depth $2$ or higher
would produce $\zet^{\mot}(2)$ or higher MZV content (Brown 2012,
iterated integral representation of motivic periods). The shadow
extraction is a \emph{residue}, not an iterated integral: residues
of simple poles at configuration boundaries give the
$(z_i \to z_j)$ OPE coefficient, not $\int \omega_{ij}$. Accordingly
no iterated-integral weight enters the extraction: the output is
depth-$1$ pure.

\emph{Motivic conclusion.} Under the stated depth-$1$ hypothesis, the
residue output lies in
$\MZV^{\mot}_0 \otimes \mathbb{Q}(\mathrm{param})$ by Brown 2012
Thm.\ 1.1 ($\MZV^{\mot}_0 = \mathbb{Q}$). This proves the asserted
weight-zero statement.

\emph{Compatibility with the induction of
Theorem~\ref{thm:virasoro-motivic-rationality-all-r}.} The
master-equation Riccati recurrence
$S_r = -(1/(r c)) \sum f(j, k) \cdot j k \cdot S_j S_k$ preserves
$\mathbb{Q}(c)$: ring-closure of $\mathbb{Q}(c)$ under addition and
multiplication. The structural theorem here explains \emph{why} the
recurrence is $\mathbb{Q}(c)$-closed: the base data (kappa, $S_3$,
$S_4$) are themselves $\mathbb{Q}(c)$-rational OPE coefficients, so
the recurrence inherits rationality from the depth-$1$ residue
structure.
\end{proof}

The structural obstruction to $\MZV^{\mot}$-content on the residue
lane is the $E_1$-versus-$E_2$ depth distinction.

\begin{theorem}[\label{thm:e1-vs-e2-mzv-depth-distinction}
$E_1$-chiral residue vs $E_2$-topological iterated integral]
\ClaimStatusConditional
Let $\cA$ be a chirally Koszul algebra on $D$. Let
$Z^{\der}_{\ch}(\cA)$ denote its derived centre, an $E_2$-chiral
algebra after Deligne--Tamarkin (cf.\ the five $E_1$-chiral notions,
), and let $\Phi_{\KZ}^{\mot}$ denote the motivic Drinfeld
associator (Alekseev--Torossian 2012, Rossi--Willwacher 2014).
\begin{enumerate}
\item[(i)] The $E_1$-chiral bar side $B^{\ord}(\cA)$ extracts
shadow residues via depth-$1$ Arnold data; its motivic image lies in
$\MZV^{\mot}_0 \otimes \mathbb{Q}(\mathrm{param}) =
\mathbb{Q}(\mathrm{param})$.
\item[(ii)] The $E_2$-topological side
$Z^{\der}_{\ch}(\cA)$ receives the Kontsevich formality quantisation,
whose structure constants populate $\MZV^{\mot}_{\geq 2}$ via
$\Phi_{\KZ}^{\mot}$; specifically, the coefficient of the
Lyndon-basis word $[x, y] [y, x, y]$ in $\Phi_{\KZ}^{\mot}$ is the
motivic period $\zet^{\mot}(3)$.
\item[(iii)] The comparison arrow $B^{\ord}(\cA) \to Z^{\der}_{\ch}
(\cA)$ \textup{(}bar to derived centre, $E_1 \to E_2$\textup{)} is
associator-dependent at the cochain level. A shadow coefficient
inherits rationality from the bar side only when it factors through
the depth-$1$ residue lane of
Theorem~\ref{thm:shadow-tower-depth-1-rationality}.
\end{enumerate}
\end{theorem}

\begin{proof}
(i) is
Theorem~\ref{thm:shadow-tower-depth-1-rationality}.

(ii) is the content of Alekseev--Torossian 2012: the Kontsevich
graph integral for the universal $L_{\infty}$-quasi-isomorphism
$\cU \colon T_{\poly} \to D_{\poly}$ produces weight-$\geq 2$
motivic periods. The wheel graph of depth $2$ contributes
$\zet^{\mot}(3)$ at the coefficient of $[x, y] [y, x, y]$; higher
wheels populate $\MZV^{\mot}_{\geq 3}$.

(iii) is the scope condition. A chiral quantum-group equivalence on
the Koszul locus may be associator-dependent on chains. Therefore a
bar-to-centre comparison cannot by itself kill positive motivic
weights. It kills them only after the coefficient is known to be
represented by the depth-$1$ residue expression of
Theorem~\ref{thm:shadow-tower-depth-1-rationality}.
\end{proof}

\begin{remark}[\label{rem:coproduct-carries-mzv-bar-does-not}
Coproduct carries MZV; bar does not]
The $E_1$-chiral bar $B^{\ord}(\cA)$ is built from deconcatenation
$T^c$ and the bar differential coming from OPE residues; no
Drinfeld associator enters its construction. The $E_1$-chiral
\emph{coproduct} $\Delta^{\ch}_z$ on $\cA$ itself (the quantum-group
datum, def:chiral-yangian-datum) does carry associator dependence in
chain-level presentations: the Etingof--Kazhdan quantisation of the
chiral coproduct uses $\Phi_{\KZ}^{\mot}$ to intertwine the two
natural $E_2$-structures on $(\cA, \cA)$. Accordingly the chiral
coproduct carries $\MZV^{\mot}$-content, while the chiral bar does
not. This is the structural source of the paradox: the same
algebra's bar-side and coproduct-side have opposite motivic
weight signatures.
\end{remark}

%%%%%%%%%%%%%%%%%%%%%%%%%%%%%%%%%%%%%%%%%%%%%%%%%%%%%%%%%%%%%%%%%%%%%%%%%%%%%
\section{Principal $\cW_N$ motivic rationality}
\label{sec:wn-motivic-rationality}
%%%%%%%%%%%%%%%%%%%%%%%%%%%%%%%%%%%%%%%%%%%%%%%%%%%%%%%%%%%%%%%%%%%%%%%%%%%%%

For the principal $\cW_N$ algebra, the shadow tower decomposes
formally into generator lines. Each generator $W_j$ of conformal
weight $j$ ($2 \leq j \leq N$) contributes a line in the
$(N-1)$-dimensional deformation space, with curvature
$\kappa_j(\cW_N) = c/j$ in the compute normalisation of
\texttt{higher\_w\_shadows.py}. Rationality on a line follows only
after that line has been placed on the depth-$1$ residue lane of
Theorem~\ref{thm:shadow-tower-depth-1-rationality}.

\begin{theorem}[\label{thm:w-n-motivic-rationality-all-r}Principal
$\cW_N$ depth-$1$ motivic rationality]
\ClaimStatusConditional
For every $N \geq 2$ and every $r \geq 2$, assume that the principal
$\cW_N$ shadow coefficient on each generator line is represented by
the depth-$1$ residue lane and that the between-line mixing terms are
rational functions of the Fateev--Lukyanov OPE constants. Then the
generator-line coefficients lie in
$\mathbb{Q}(c)$:
\[
 S_r^{\cW_N, T}(c), \; S_r^{\cW_N, W_3}(c), \; \ldots, \;
 S_r^{\cW_N, W_N}(c) \;\in\; \mathbb{Q}(c),
\]
On this conditional lane no $\zet^{\mot}$ or $(2\pi i)^{\mot}$
content appears.
\end{theorem}

\begin{proof}
Each generator line $W_j$ inherits a one-parameter sub-tower
governed by the shadow data $(\kappa_j, \alpha_j, S_4^{W_j})$ from
\texttt{higher\_w\_shadows.py}. These data are $\mathbb{Q}(c)$-rational
by the Fateev--Lukyanov $\cW_N$ OPE: structure constants are rational
in $c$ (Fateev--Lukyanov 1988, Appendix A). The shadow recurrence on
a single line reduces to the Riccati master equation at the line's
effective kappa $\kappa_j = c/j$, with base data in $\mathbb{Q}(c)$.
Under the stated depth-$1$ hypothesis,
Theorem~\ref{thm:shadow-tower-depth-1-rationality} gives
$S_r^{\cW_N, W_j}(c) \in \mathbb{Q}(c)$ and kills positive motivic
weight on that line. The Virasoro induction on the Riccati master
equation supplies the model calculation; it becomes a proof for
$W_j$ only after the $W_j$ residue has been identified with the same
depth-$1$ expression.

\emph{Between-line mixing.} Off-line shadow data involves the
propagator variance $\delta_{\mathrm{mix}}$. Rationality of this term
is part of the stated depth-$1$ mixing hypothesis; it is not a formal
consequence of Cauchy--Schwarz.

\emph{Depth-$1$ structural guarantee.} Under the stated hypothesis,
the residue extraction on $\mathrm{Conf}_n(D)$ is depth-$1$ pure for
the chosen generator line. Theorem~\ref{thm:shadow-tower-depth-1-rationality}
then gives $\mathbb{Q}(c)$-rationality on that line.
\end{proof}

\begin{proposition}[\label{prop:w3-w-line-motivic-rationality}
$\cW_3$ W-line explicit rationality]
\ClaimStatusProvedHere
The $\cW_3$ W-line shadow coefficients vanish at odd $r$ (by
$\mathbb{Z}_2$ parity $W \mapsto -W$) and at even $r = 2n \geq 4$
admit the closed form
\[
 S_{2n}^{\cW_3, W}(c) \;=\; (-1)^n a_n \cdot \frac{2560^{n-1}}
 {c^{2n-3} (5c + 22)^{3(n-1)}}
 \;\in\; \mathbb{Q}(c),
\]
with $a_n = \tfrac{(-12)^n}{54} \cdot \binom{3/2}{n}$. The denominator
divides a power of $c (5c+22)^3$ at every weight, with no new
irreducible factor entering.
\end{proposition}

\begin{proof}
Reading from \texttt{compute/lib/w3\_shadow\_tower\_arity12\_engine.py}
(function \texttt{w\_line\_general\_term}), the even-$r$ W-line
admits the explicit formula above. Direct substitution gives
$S_4 = 2560/[c(5c+22)^3]$, $S_6 = -2 S_4^2/c$, $S_8 = 9 S_4^3/c^2$,
$S_{10} = -54 S_4^4/c^3$, $S_{12} = 378 S_4^5/c^4$ (ring relations
verified in \texttt{w\_line\_ring\_coefficients}). Each entry is
$\mathbb{Q}(c)$-rational; the denominator contains only $c$ and
$(5c+22)$. Odd-$r$ vanishing is the $\mathbb{Z}_2$ parity reflection
$W \mapsto -W$ of the $\cW_3$ OPE at odd $W$-insertion count.
\end{proof}

The W-line case delivers an additional structural insight: no MZV
emergence at any probed $r$ through $r = 12$ in the direct
closed-form expansion. Depth-$2$ Arnold integrals would manifest as
$\zet^{\mot}(n)$ atoms; none arise.

%%%%%%%%%%%%%%%%%%%%%%%%%%%%%%%%%%%%%%%%%%%%%%%%%%%%%%%%%%%%%%%%%%%%%%%%%%%%%
\section{Bershadsky--Polyakov motivic rationality}
\label{sec:bp-motivic-rationality}
%%%%%%%%%%%%%%%%%%%%%%%%%%%%%%%%%%%%%%%%%%%%%%%%%%%%%%%%%%%%%%%%%%%%%%%%%%%%%

The Bershadsky--Polyakov algebra $\mathrm{BP}_k = \cW^k(sl_3,
f_{(2,1)})$ is the minimal-nilpotent DS reduction of $\widehat{sl_3}$;
its four generators $(J, G^+, G^-, T)$ have weights $(1, 3/2, 3/2, 2)$
and its central charge in Arakawa convention is
$c_{\mathrm{BP}}(k) = 2 - 24(k+1)^2/(k+3)$
(Fehily--Kawasetsu--Ridout 2020). The Koszul conductor
$K_{\mathrm{BP}}(k) = c_{\mathrm{BP}}(k) + c_{\mathrm{BP}}(-k-6) = 196$
is level-independent
(\texttt{bp\_self\_duality.tex} Thm.\
\texttt{thm:bp-koszul-conductor-polynomial}).

\begin{theorem}[\label{thm:bp-motivic-rationality-arakawa}BP T-line
motivic rationality in Arakawa convention]
\ClaimStatusProvedHere
For every $r \geq 2$, the BP shadow coefficient on the T-line
(Virasoro restriction) lies in $\mathbb{Q}(k)$:
\[
 S_r^{\mathrm{BP}, T}(k) \;=\;
 S_r^{\Vir}\!\bigl(c_{\mathrm{BP}}(k)\bigr)
 \;\in\;
 \mathbb{Q}(k),
\]
with motivic lift of weight zero and no $\zet^{\mot}$ content at any
weight. The denominator factors only through
$(k+3)$ and a finite family of rational expressions in
$c_{\mathrm{BP}}(k)$.
\end{theorem}

\begin{proof}
The T-line shadow of BP reduces by restriction to the Virasoro
shadow at the BP central charge: $S_r^{\mathrm{BP}, T}(k) =
S_r^{\Vir}(c_{\mathrm{BP}}(k))$ (\texttt{bp\_shadow\_tower.py}
function \texttt{bp\_tline\_shadow\_tower}). Apply
Theorem~\ref{thm:virasoro-motivic-rationality-all-r}: at any
$\mathbb{Q}$-rational specialisation $c = c_{\mathrm{BP}}(k)$ of
$c \in \mathbb{Q}(c)$, the output $S_r^{\Vir}(c_{\mathrm{BP}}(k))$
remains $\mathbb{Q}$-rational in $k$ because substitution of a
$\mathbb{Q}(k)$-rational central charge into a $\mathbb{Q}(c)$-rational
function produces a $\mathbb{Q}(k)$-rational function. The motivic
primitivity is preserved under base change of the parameter field
$\mathbb{Q}(c) \to \mathbb{Q}(k)$ via the substitution $c \mapsto
c_{\mathrm{BP}}(k)$; Brown's weight-grading is parameter-field
independent.

The theorem concerns only the T-line. The $J$ and $G^\pm$ lines require
their own residue-lane identifications before a class-M rationality
statement can be made for them.
\end{proof}

\begin{proposition}[\label{prop:bp-fl-convention-caveat}FL-convention
Koszul conductor: distinct constant]
\ClaimStatusProvedHere
Under the Fateev--Lukyanov screening convention
$c_{\mathrm{BP}}^{\mathrm{FL}}(k) = -(2k+3)(3k+1)/(k+3)$, the
Koszul-conductor identity under the level involution $k \mapsto -k-6$
is the polynomial identity
\[
 K_{\mathrm{BP}}^{\mathrm{FL}}(k) \;=\;
 c_{\mathrm{BP}}^{\mathrm{FL}}(k) + c_{\mathrm{BP}}^{\mathrm{FL}}(-k-6)
 \;=\; 50
 \qquad \text{(constant in } k \text{).}
\]
Both conventions parametrise the \emph{same} algebra via level
reparametrisation; Theorem~\ref{thm:bp-motivic-rationality-arakawa}
therefore descends to the FL convention through the rational level
reparametrisation, and the Koszul-conductor identity remains
polynomial in both. The \emph{value} of the conductor depends on
convention: $196$ in Arakawa, $50$ in Fateev--Lukyanov; this
convention shift reflects a redefinition of the central charge and
the scale of the dual-level pairing, not a change in the algebra.
\end{proposition}

\begin{proof}
Direct substitution: writing
$c_{\mathrm{BP}}^{\mathrm{FL}}(k) = (-6k^2 - 11k - 3)/(k + 3)$ and
$c_{\mathrm{BP}}^{\mathrm{FL}}(-k-6) = (6k^2 + 61k + 153)/(k + 3)$,
their sum is $(50(k+3))/(k+3) = 50$ after cancellation of the common
denominator. This is a polynomial identity in $k$ (in fact constant),
verified in the compute engine
\texttt{compute/tests/test\_motivic\_shadow\_full\_class\_m.py} at
four explicit levels $k \in \{1, 2, 4, 10\}$ and via the
\texttt{sympy.simplify} symbolic identity. Volume~II's
\texttt{bp\_chain\_level\_strict\_platonic.tex} uses the Fateev--Lukyanov
convention consistently. Rational change of level variable preserves
$\mathbb{Q}(k)$-membership, so the motivic-rationality conclusion of
Theorem~\ref{thm:bp-motivic-rationality-arakawa} descends.
\end{proof}

\begin{remark}[\label{rem:bp-conductor-polynomial-constant}Polynomial-constant
identity in both conventions]
Direct substitution of $c^{\mathrm{FL}}_{\mathrm{BP}}(k)$ and
$c^{\mathrm{FL}}_{\mathrm{BP}}(-k-6)$ yields the polynomial-constant
value $K^{\mathrm{FL}}_{\mathrm{BP}} = 50$: any putative meromorphic
expression $-12(k+3) - 48/(k+3)$ is an arithmetic artefact, as the
apparent pole at $k = -3$ cancels in the sum.
Proposition~\ref{prop:bp-fl-convention-caveat} records the polynomial
identity $K^{\mathrm{FL}}_{\mathrm{BP}} = 50$. The Arakawa convention
gives $K^{\mathrm{Arakawa}}_{\mathrm{BP}} = 196$; the two constants
differ by a convention-dependent rescaling ($196 - 50 = 146 =
2 \cdot 73$, the factor accompanying the change of central-charge
normalisation between FL screening and Arakawa DS reduction).
\end{remark}

%%%%%%%%%%%%%%%%%%%%%%%%%%%%%%%%%%%%%%%%%%%%%%%%%%%%%%%%%%%%%%%%%%%%%%%%%%%%%
\section{$\cW_{\infty}$ motivic rationality at generic Miura parameter}
\label{sec:w-infty-motivic-rationality}
%%%%%%%%%%%%%%%%%%%%%%%%%%%%%%%%%%%%%%%%%%%%%%%%%%%%%%%%%%%%%%%%%%%%%%%%%%%%%

The $\cW_{\infty}$ algebra at Miura parameter $\Psi$ admits a
triality decomposition; its shadow tower is governed by the Miura
cross-universality identity
$\Delta_z(T) = T \otimes 1 + 1 \otimes T + ((\Psi - 1)/\Psi) \cdot J
\otimes J$ (Vol~I \texttt{miura\_cross\_universality.tex}
\texttt{thm:miura-cross-universality}).

\begin{theorem}[\label{thm:w-infty-motivic-rationality-all-r}
$\cW_{\infty}$ depth-$1$ motivic rationality]
\ClaimStatusConditional
For every $r \geq 2$ and at generic Miura parameter $\Psi$, assume
that the large-$N$ limit of the principal $\cW_N$ depth-$1$ residue
lane exists coefficientwise and that Miura cross-terms remain rational
in $(c,\Psi)$. Then the
$\cW_{\infty}$ shadow coefficient on the T-line satisfies
\[
 S_r^{\cW_{\infty}, T}(c, \Psi) \;\in\; \mathbb{Q}(c, \Psi),
\]
with no $\zet^{\mot}$ content on this lane. The Miura parameter
$\Psi$ enters only through
the rational prefactor $(\Psi - 1)/\Psi$ of the cross-arity
Stasheff identity, not through any MZV-weighted period.
\end{theorem}

\begin{proof}
Under the stated hypothesis, the $\cW_{\infty}$ shadow coefficient on
the T-line is the coefficientwise large-$N$ limit of the $\cW_N$
T-line coefficient (\texttt{w\_infinity\_shadow\_limit.py}); the
conditional Theorem~\ref{thm:w-n-motivic-rationality-all-r} then
places the limit in $\mathbb{Q}(c)$. The Miura cross-universality of
\texttt{thm:miura-cross-universality} adds the rational prefactor
$(\Psi - 1)/\Psi$ at every spin and cross-arity; this is a
$\mathbb{Q}(\Psi)$-rational function of $\Psi$ with no
$\MZV^{\mot}$-content. The combined $\mathbb{Q}(c) \cdot
\mathbb{Q}(\Psi)$-structure gives $\mathbb{Q}(c, \Psi)$.

\emph{Generic-$\Psi$ scope.} At $\Psi = 0$ the algebra degenerates
to the free abelian chiral algebra with identity coproduct
(Vol~I \texttt{concordance.tex} entry on $\Psi = 0$); at $\Psi = 1$
the free-boson point gives the OPE-obstruction-vanishing free-field
realisation. At these special values the shadow tower trivialises
($S_r^{\cW_{\infty}, T}(c, 0) = 0$ for $r \geq 3$; at $\Psi = 1$, the
shadow becomes the Heisenberg class G tower). At generic $\Psi$
(neither $0$ nor $1$) the tower is class M and inherits rationality
only on the coefficientwise depth-$1$ large-$N$ lane.

\emph{Depth-$1$ structure.} The Miura screening intersection
produces $\cW_{\infty}$ as a sub-algebra of the rank-$\infty$
Heisenberg; the OPE data on the Miura presentation are rational in
$(c, \Psi)$. This rationality is not by itself the theorem. The
remaining hypothesis is precisely that the coefficient under
discussion is computed by the depth-$1$ residue expression of
Theorem~\ref{thm:shadow-tower-depth-1-rationality}.
\end{proof}

\begin{remark}[\label{rem:w-infty-psi-specializations}Special-$\Psi$
specialisations]
At $\Psi = 2$ (free-fermion point) the $\cW_{\infty}$ shadow reduces
to a symplectic-fermion tower, class C, terminating at $r = 4$. At
$\Psi$ equal to an integer $\geq 3$ with $(\Psi, c)$ in the
admissible locus of Arakawa, the shadow reduces to the admissible
finite-length KL block. Motivic rationality persists only for those
specialised coefficients which remain on the depth-$1$ residue lane.
These specialisations lie within the $\mathbb{Q}(c,\Psi)$-rational
locus of Theorem~\ref{thm:w-infty-motivic-rationality-all-r} after
that lane condition is checked.
\end{remark}

%%%%%%%%%%%%%%%%%%%%%%%%%%%%%%%%%%%%%%%%%%%%%%%%%%%%%%%%%%%%%%%%%%%%%%%%%%%%%
\section{Conditional class-M depth-$1$ rationality}
\label{sec:class-m-motivic-rationality-full}
%%%%%%%%%%%%%%%%%%%%%%%%%%%%%%%%%%%%%%%%%%%%%%%%%%%%%%%%%%%%%%%%%%%%%%%%%%%%%

The family statements combine only after the residue-lane hypothesis is
kept in the statement.

\begin{theorem}[\label{thm:class-m-motivic-rationality-full}
Class M depth-$1$ motivic rationality]
\ClaimStatusConditional
Let $\cA$ be a chirally Koszul algebra in class M with
$\mathbb{Q}$-rational OPE structure constants. Fix $r \geq 2$ and
assume that $S_r(\cA)$ is represented on the depth-$1$ residue lane of
Theorem~\ref{thm:shadow-tower-depth-1-rationality}. Then
$S_r(\cA)$ lies in $\mathbb{Q}(\mathrm{param})$ and its motivic lift
has no positive-weight MZV component.
\end{theorem}

\begin{proof}
This is Theorem~\ref{thm:shadow-tower-depth-1-rationality} applied
with the stated hypothesis. The Virasoro and BP T-line cases give
proved examples of such a residue lane; the principal $\cW_N$ and
$\cW_\infty$ statements above record the conditional extension to the
generator-line and Miura-limit surfaces. No assertion is made that
every class-M coefficient is depth-$1$.
\end{proof}

%%%%%%%%%%%%%%%%%%%%%%%%%%%%%%%%%%%%%%%%%%%%%%%%%%%%%%%%%%%%%%%%%%%%%%%%%%%%%
\section{Tempered Virasoro and the motivic frontier}
\label{sec:tempered-vir-motivic-frontier}
%%%%%%%%%%%%%%%%%%%%%%%%%%%%%%%%%%%%%%%%%%%%%%%%%%%%%%%%%%%%%%%%%%%%%%%%%%%%%

\begin{remark}[\label{rem:tempered-virasoro-not-realized}Tempered
Virasoro is not an independent class-M algebra]
``Tempered Virasoro'' is not an independent standard-landscape
member. Candidate constructions (Virasoro at tempered representations
in the sense of Harish--Chandra, or Virasoro at the $c \to \infty$
limit with $c$-regularised shadow) reduce in all cases to
Virasoro at a reparametrised central charge. By
Theorem~\ref{thm:virasoro-motivic-rationality-all-r} the shadow
tower stays $\mathbb{Q}(c)$-rational under any such
reparametrisation, so tempered Virasoro (if one constructs it) would
inherit motivic rationality automatically; no new theorem is needed.
\end{remark}

\begin{remark}[\label{rem:motivic-frontier-class-m}Motivic frontier]
The structural Theorem~\ref{thm:shadow-tower-depth-1-rationality}
applies to coefficients already known to lie on the depth-$1$ residue
lane. Two classes escape the easy rationality mechanism:
\begin{enumerate}
\item[(a)] \emph{Transcendentally-parameterised algebras.} If the
OPE structure constants involve $\pi$, $\zet(3)$, or other
transcendental periods (e.g.\ a Moonshine-bimodule algebra at
$\pi$-dependent twists), then the shadow tower can carry
$\MZV^{\mot}$-content. No such algebra appears in the standard
landscape.
\item[(b)] \emph{Non-compact algebras.} At $c \to \infty$ in
Virasoro, the regularised shadow $S_r(\Vir_c)/c^{r/2}$ stabilises
(see \texttt{virasoro\_c\_to\_infty.py}); the limit is
$\mathbb{Q}$-rational, so rationality is preserved.
\end{enumerate}
The genuine frontier is: do there exist chirally Koszul algebras
with irrational-algebraic (but not transcendental) OPE structure
constants, and if so does the shadow tower see them as
$\mathbb{Q}$-algebraic or $\MZV^{\mot}$-valued? We conjecture
algebraic-irrational OPE constants (e.g.\ $\sqrt{2}$-rational
structure constants from an $sl_2$-type free-field construction)
produce $\mathbb{Q}(\sqrt{2}, \mathrm{param})$-rational shadow
coefficients: algebraic but still depth-$1$ pure.
\end{remark}

%%%%%%%%%%%%%%%%%%%%%%%%%%%%%%%%%%%%%%%%%%%%%%%%%%%%%%%%%%%%%%%%%%%%%%%%%%%%%
\section{Cross-volume coherence}
\label{sec:cross-volume-coherence}
%%%%%%%%%%%%%%%%%%%%%%%%%%%%%%%%%%%%%%%%%%%%%%%%%%%%%%%%%%%%%%%%%%%%%%%%%%%%%

The conditional class-M depth-$1$ theorem interacts with Vol~II's
gravitational climax and Vol~III's BPS landscape at the following
junctions.

\begin{enumerate}
\item[(i)] \emph{Vol~II gravitational climax.} The 3d HT topological
theory carries motivic periods of the Chern--Simons partition
function at integer level; these are computed from the $E_2$
topological side (Kontsevich formality on the boundary
$\cW_{\infty}$-algebra), and populate $\MZV^{\mot}_{\geq 2}$. The
shadow-tower component stays $\mathbb{Q}$-rational only on the
depth-$1$ lane of
Theorem~\ref{thm:class-m-motivic-rationality-full}. The division of
motivic labour is conditional: residue-bar coefficients are rational,
while coproduct and higher-period coefficients may carry
$\MZV^{\mot}$.
\item[(ii)] \emph{Vol~III BPS landscape.} CY$_3$ BPS algebra
shadow on the Koszul locus matches Vol~I's shadow tower via the
Vol~III two-stage factorisation
$\Phi_d^{(\Sigma_{d-1},C)} = \SpCh_{\Sigma_{d-1}, C} \circ \PhiFA_d$
(Stage~1 canonical $E_d$-holomorphic factorisation algebra on the
CY$_d$ variety, Stage~2 specialisation to a reference curve~$C$
via factorisation homology over a $(d{-}1)$-cycle); the Vol~III
$\kappa_{\ch}$
is the Vol~I $\kappa$ under this identification, and motivic
rationality propagates only for the matched depth-$1$ residue
coefficients. BKM lattice shadow coefficients carry
$(2\pi i)^{\mot, 2}$ weights via Borcherds's theta-lift, but these
are coproduct-tower contributions (category-level $\kappa_{\mathrm{cat}}$),
not bar-tower contributions.
\end{enumerate}

\begin{remark}[\label{rem:resolves-item-19-full-class-m}Conditional
resolution of item 19 on the depth-$1$ lane]
The motivic frontier asks for MZV-freeness
of the shadow tower. It is closed only on the depth-$1$ residue lane:
Theorems~\ref{thm:virasoro-motivic-rationality-all-r},
\ref{thm:w-n-motivic-rationality-all-r},
\ref{thm:bp-motivic-rationality-arakawa},
\ref{thm:w-infty-motivic-rationality-all-r}, and
\ref{thm:class-m-motivic-rationality-full} cover the stated residue
surfaces, and Theorem~\ref{thm:shadow-tower-depth-1-rationality}
supplies the structural mechanism. The $E_1$-vs-$E_2$ distinction of
Theorem~\ref{thm:e1-vs-e2-mzv-depth-distinction} explains the
paradox only after the depth of the coefficient has been fixed:
motivic richness can live on the coproduct side or on higher
configuration-space periods, while the proved rationality concerns
the residue-bar lane.
\end{remark}
